\BourbakiExerciseGroup

\begin{enumerate}
\def\labelenumi{\arabic{enumi})}
\tightlist
\item
  设 X 是无限集，F 是 X 上的一个超滤子，且 F 中诸集之交为空。对每个 \(A \in F\)，设 \(V_{A}\) 表示子集 \(\Delta \cup (A \times A)\)（\(X \times X\) 的）。证明：当 A 取遍 F 时，诸集合 \(V_{A}\) 构成 X 上某个一致结构 \(\mathcal{U}(\mathfrak{F})\) 的一致邻域基本系，且该一致结构诱导的拓扑是离散拓扑。
\end{enumerate}

*2) 在带有加法一致结构的实直线 R 上，证明当 V 取遍 R 的一致邻域滤子时，诸集合 V(Z) 并不构成有理整数集 Z 的基本邻域系。*（参看 § 4 第 3 目命题 4 的推论）。

*3) 设 V 是 R 上加法一致结构中由所有对 \((x, y)\)（满足 \(|x - y| \leqslant 1\) 或 \(xy \geqslant 1\) 的）组成的一致邻域。证明 V 在 \(R \times R\) 中是闭的，但若 A 表示由所有整数 \(n \geqslant 2\) 组成的 R 的（闭）子集，则 \(V(A)\) 在 \(R\) 中不是闭的。*

\begin{enumerate}
\def\labelenumi{\arabic{enumi})}
\setcounter{enumi}{3}
\tightlist
\item
  证明若一致空间是 Kolmogoroff 空间（第一章 § 1 习题 2），则它是 Hausdorff 的。
\end{enumerate}

¶ 5) a) 设 X 是一致空间，U 是它的一致结构；对 X 的每个一致邻域 V，设 \(\tilde{V}\) 是 \(\mathfrak{P}(X) \times \mathfrak{P}(X)\) 的子集，由 X 的子集的所有满足 \(M \subset V(N)\) 且 \(N \subset V(M)\) 的对 (M, N) 组成。证明诸集合 \(\tilde{V}\) 构成某个一致结构 \(\tilde{U}\)（\(\mathfrak{P}(X)\) 上的）的一致邻域基本系。

\begin{enumerate}
\def\labelenumi{\alph{enumi})}
\setcounter{enumi}{1}
\item
  在集 \(\mathfrak{P}_{0}(\mathbf{X})\)（由 \(\mathbf{X}\) 的非空子集组成）上，证明由拓扑 \(\mathfrak{C}(\tilde{\mathcal{U}})\)（它由 \(\tilde{\mathcal{U}}\) 诱导）所诱导的拓扑比拓扑 \(\mathfrak{C}_{\Phi}\) 细（第一章 § 2 习题 7）。
\item
  若 X 至少有两点且 U 是 Hausdorff 的，证明 \(\mathfrak{G}(\tilde{\mathcal{U}})\) 在 \(\mathfrak{P}_{0}(\mathbf{X})\) 上诱导的拓扑决不比 \(G_{\Omega}\) 粗。\(\mathfrak{G}(\tilde{\mathcal{U}})\) 在由 X 的非空闭子集组成的集 \(\mathfrak{F}(\mathbf{X})\) 上诱导的拓扑比 \(G_{\Omega}\) 在这个集上诱导的拓扑细，当且仅当对 X 的每个闭子集 A，当 V 取遍 X 的一致邻域时诸集合 V(A) 构成 A 的一个基本邻域系。
\end{enumerate}

*d) 证明在第一章 § 11 习题 18 定义的商集 X/R 上，商拓扑以及由 \(\mathfrak{C}_{\Omega}\)、\(\mathfrak{C}_{\Phi}\) 与 \(\mathfrak{C}(\tilde{\mathcal{U}})\)（其中 \(\mathcal{U}\) 是 \(\mathbf{R}^2\) 上的通常一致结构）定义的拓扑，这四个拓扑互不相同。*

